% cover.tex - title page and contents of the combined PDF
\pagenumbering{roman}
\begin{titlepage}
\centering
\vspace*{3.5cm}
{\color{brand}\rule{\linewidth}{1.2pt}}\par\vspace{16pt}
{\sffamily\bfseries\fontsize{30}{38}\selectfont\color{brand}Raven--MODFLOW~6\\[4pt]Tutorials\par}
\vspace{14pt}
{\color{brand}\rule{\linewidth}{1.2pt}}\par\vspace{24pt}
{\large Six click-by-click tutorials for the groundwater coupling of the\\[2pt]Raven Hydrological Modelling Framework\par}
\vfill
{\large November 2026\par}
\vspace{1.2cm}
{\footnotesize\color{gray!70!black}Licensed under the Creative Commons Attribution 4.0 International License (CC BY 4.0),
\url{https://creativecommons.org/licenses/by/4.0/}.\par}
\end{titlepage}
\renewcommand{\tutheader}{Raven--MODFLOW~6 tutorials}
\begingroup\small\tableofcontents\endgroup
\chapter*{About these tutorials}
\addcontentsline{toc}{chapter}{About these tutorials}
Each tutorial is a sequence of steps typed in a Linux terminal (Windows: the Windows Subsystem for Linux). Every file you
create or change is shown in full or as an exact difference, every command with what it prints, and every number and
figure comes from running the steps with the Raven of the release package; a rerun of \cmd{build\_tutorials.py}
regenerates them.
\begin{enumerate}
\item \textbf{Getting started:} install the package, check it, run the Liard River example and read its groundwater outputs.
\item \textbf{Your own groundwater model:} give a Raven model (Tutorial Creek) a groundwater model that Raven builds.
\item \textbf{The soil-zone coupling:} let the water table rise into the soil; unsaturated columns below it.
\item \textbf{Coupling an existing MODFLOW~6 model:} link the HRUs to a model in feet; ET, a wetland and a reservoir.
\item \textbf{Restarts, ensembles and particle tracking:} continue runs exactly, run a Monte Carlo ensemble, track particles.
\item \textbf{Calibration with Ostrich (DDS):} recover known aquifer parameters from streamflow and heads.
\end{enumerate}
Tutorial~1 comes first; Tutorials 3 to 6 start from the model built in Tutorial~2, whose finished state is included, so
any of them can be done on its own. The tutorial files are in \cmd{Raven\_MF6\_Tutorial\_Files.zip}: one folder per
tutorial with its \cmd{start} and \cmd{finished} state, and the scripts the tutorials use.
\cleardoublepage
\pagenumbering{arabic}
